\begin{problem}[A nonrational closed point]\label{prob:vi17-p15}
For $A=\mathbb R[x]$ and $\mfm=(x^2+1)$, compute $\OO_{X,\mfm}$ and $\kappa(\mfm)$, and interpret the value of a real polynomial at this scheme point.
\end{problem}

\ifdefined\FullProblemDossiers
\paragraph{Why this problem matters.}
This problem breaks the habit of imagining every closed point as taking values in the ground field.

\paragraph{Useful ideas.}
Localize first, then quotient by the maximal ideal.

\paragraph{Guided hints.}
\begin{enumerate}
\item The stalk is $\mathbb R[x]_{(x^2+1)}$.
\item The residue field is the quotient by $x^2+1$.
\end{enumerate}

\begin{solution}
The stalk theorem gives $\OO_{X,\mfm}=\mathbb R[x]_{(x^2+1)}$.  Its residue field is
\[\mathbb R[x]_{(x^2+1)}/(x^2+1)\cong\mathbb R[x]/(x^2+1)\cong\mathbb C.\]
Thus the value of $p(x)$ at this point is the class of $p(x)$ in $\mathbb C$, equivalently $p(i)$ after choosing the usual identification.
\end{solution}

\paragraph{What happened mathematically.}
Scheme points naturally carry extension fields, so evaluation is residue-field valued.

\paragraph{Extensions.}
Do the same for a maximal ideal generated by an irreducible polynomial over an arbitrary field.

\paragraph{Related problems.}
VI.17.P14 and VI.17.P16 develop neighboring aspects of the same localization-sheaf calculus.

\paragraph{Provenance.}
\textbf{New synthesis of a recurring corpus residue-field example.}
\else
\paragraph{Provenance.} \textbf{New synthesis of a recurring corpus residue-field example.}
\fi
